%%%%%%%%%%%%%%%%%%%%%%%%%%%%%%%%%%%%%%%%%%%%%%%%%%%%%%%%%%
%%
%% MARK: Title & Data
%%
%%%%%%%%%%%%%%%%%%%%%%%%%%%%%%%%%%%%%%%%%%%%%%%%%%%%%%%%%%

\addtocontents{toc}{\SkipTocEntry}
\chapter*{\S B34. Smooth Embeddings and \hbox{Smoothly Embedded Subsets}}
\chaptermark{Smooth embeddings and smoothly embedded subsets}
\addcontentsline{toc}{chapter}{Smooth embeddings and smoothly embedded subsets}
\setcounter{footnote}{0}

%%%%%%%%%%%%%%%%%%%%%%%%%%%%%%%%%%%%%%%%%%%%%%%%%%%%%%%%%%
%%
%% MARK: Abstract
%%
%%%%%%%%%%%%%%%%%%%%%%%%%%%%%%%%%%%%%%%%%%%%%%%%%%%%%%%%%%

\begin{abstract}
  Since Henri Joris's paper we have known that,
  paradoxically,
  the semi-cubic $y^2 - x^3 = 0$ is an embedded submanifold of $\RR^2$.
  So,
  how are we to understand the singularity of the cusp?
  The answer is that,
  although the cusp is embedded,
  it is not smoothly embedded.
  This note will detail this aspect.
\end{abstract}

%%%%%%%%%%%%%%%%%%%%%%%%%%%%%%%%%%%%%%%%%%%%%%%%%%%%%%%%%%
%%
%% MARK: Smooth Embeddings
%%
%%%%%%%%%%%%%%%%%%%%%%%%%%%%%%%%%%%%%%%%%%%%%%%%%%%%%%%%%%

\section{Smooth embeddings}

The semi-cubic is the subset $\fS \subset \RR^2$ defined by the equation $y^2 - x^3 = 0$.
It is represented in Figure \ref{fig-semi-cubic}.
It is the image of the map
$$
  j : \RR \to \RR^2, \quad j(t) = \Vector{t^2 \\ t^3}.
  $$
The map $j$ is not an immersion since $j'(0) = 0$,
where $j'$ denotes the derivative of $j$.
But a long-standing question is:
is $j$ an induction?

Recall that a map $j$ is an induction \cite[\textsection\,1.29]{PIZ13} if and only if
$j$ is smooth,
injective,
and the inverse $j^{-1} : \fS \to \RR$ is smooth,
where $\fS = \Val(j)$ is equipped with the subset diffeology.

In this case,
this is equivalent to the following:

\Question{Question.} For any smooth parametrization $P : r \mapsto (x(r),y(r))$ with values in $\fS$,
that is,
$y(r)^2 - x(r)^3 = 0$ for all $r$,
is the following map smooth?
$$
  t(r) = {y(r) \over x(r)}.
  $$
The solution to this question comes from a theorem due to Henri Joris \cite{Jor82},
which was pointed out to me by Yael Karshon.%
\footnote{In a private discussion,
but since commented on in \cite{KMW22}.}
In a simple version:

\begin{Article}{Theorem (H. Joris).}{\itshape
  If a real map $f$ is such that $t \mapsto f(t)^2$ and $t \mapsto f(t)^3$ are smooth,
  then $f$ is smooth.
  }
\end{Article}

\begin{figure}[htb]
  \centerline{\includegraphics[width=0.33\textwidth]{Figures/fig-semi-cubic}}
  \caption{The semi-cubic $y^2 = x^3$.}
  \label{fig-semi-cubic}
\end{figure}

\begin{Article}{Corollary.}{\itshape
  The map $j : t \mapsto (t^2,t^3)$ is an induction from $\RR$ into $\RR^2$.
  The semi-cubic $\fS$, the image of $j$, is an embedded submanifold of $\RR^2$.
  }
\end{Article}

\begin{proof}
  Clearly $j : t \mapsto (x = t^2, y = t^3)$ is smooth and injective:
  $t = \sqrt[3]{y}$.
  Let $r \mapsto P(r) = (x(r),y(r))$ be a plot in $\RR^2$ with values in $\fS$.
  Then $j^{-1} \circ P(r) = t(r)$ is such that $r \mapsto t(r)^2$ and $r \mapsto t(r)^3$ are smooth.
  To apply Joris's theorem we need to return to a map from $\RR$ to $\RR$.
  We can use Boman's theorem \cite{Bom67}:
  
  \begin{Article}{Theorem (J. Boman).}{\itshape
    A parametrization $P : U \to \RR^m$ is smooth if and only if,
    for every smooth path $\gamma$ in $U$,
    the composite $P \circ \gamma$ is smooth.
    }
  \end{Article}
  
  For each such smooth path,
  the composite $s \mapsto t(r(s))$ is smooth,
  thanks indeed to Joris's theorem.
  Thus $r \mapsto t(r)$ is smooth.
  Finally, $j$ is an embedding because $y \mapsto \sqrt[3]{y}$ is a homeomorphism of $\RR$.
\end{proof}

Thus,
we have a figure,
the semi-cubic,
the image of $\RR$ by an induction,
which is by construction a submanifold of $\RR^2$ because its subset diffeology is equivalent to that of $\RR$.

There is clearly something strange about this situation:
the cusp at $(0,0)$ is obviously a singularity.
But since it is transparent to the subset diffeology,
how are we to capture it?
That is the question.

%%%%%%%%%%%%%%%%%%%%%%%%%%%%%%%%%%%%%%%%%%%%%%%%%%%%%%%%%%
%%
%% MARK: Smoothly Embedded Subsets
%%
%%%%%%%%%%%%%%%%%%%%%%%%%%%%%%%%%%%%%%%%%%%%%%%%%%%%%%%%%%

\section{Smoothly embedded subsets}

The answer lies in the relationship between the ambient space $\RR^2$ and the subspace $\fS$.
Consider the pseudogroup of local diffeomorphisms of $\RR^2$ that preserve $\fS$,
that is,
that globally fix $\fS$:
$$
  \Diff_\loc(\RR^2,\fS) = \{\pphi \in \Diff_\loc(\RR^2) \mid \pphi(\fS) \subset \fS \}.
  $$
Since $\pphi$ is a local diffeomorphism, $\pphi(\fS) = \pphi(\fS \cap \Domain(\pphi))$.
If $\fS \cap \Domain(\pphi)$ is empty,
there is nothing to check.
Then,

\begin{Article}{Proposition.}{}
  Every local diffeomorphism $\pphi \in \Diff_\loc(\RR^2,\fS)$ fixes the point $(0,0)$:
  $$
    \forall \pphi \in \Diff_\loc(\RR^2,\fS),\ \pphi(0,0) = (0,0).
    $$
\end{Article}

\begin{proof}
  Let $\pphi \in \Diff_\loc(\RR^2,\fS)$ be such that $\pphi(0,0) = (x,y)$ with $(x,y) \neq (0,0)$.
  The restriction $\pphi \restriction \fS$ belongs to $\Diff_\loc(\fS)$.
  Thus $\varphi = j^{-1} \circ (\pphi \restriction \fS) \circ j$ belongs to $\Diff_\loc(\RR)$.
  \begin{center}
    \begin{tikzcd}[column sep=large, row sep=large, every label/.append style = {font = \small}]
      \RR \arrow[r,"j"]  &  \fS \subset \RR^2  \\
      \RR \arrow[u, dashed, "\varphi"] \arrow[r, swap, "j"] & \fS \subset \RR^2 \arrow[u, swap, dashed, "\pphi \restriction\fS"]
    \end{tikzcd}
  \end{center}
  Differentiate $j \circ \varphi = (\pphi \restriction \fS) \circ j$ at $0$,
  and set $j(t) = (x,y) = \pphi(0,0)$:
  \begin{eqnarray*}
    \Der(j \circ \varphi)(0) & = & \Der((\pphi \restriction \fS) \circ j)(0) \\
    \Der(j)(t) \circ \Der(\varphi)(0) & = & \Der(\pphi \restriction \fS)(0,0) \circ \Der(j)(0) \\
    \Der(j)(t) \circ \Der(\varphi)(0)  & = & 0.
  \end{eqnarray*}
  But $t \neq 0$ implies $\Rank(\Der(j)(t)) = 1$,
  and $\Der(\varphi)(0)$ is a non-zero number,
  so the composite $\Der(j)(t) \circ \Der(\varphi)(0)$ cannot be zero.
  This contradiction shows that $(x,y) = (0,0)$.
\end{proof}

Thus,
although the local diffeomorphisms are transitive on $\fS$,
the local diffeomorphisms of $\RR^2$ that preserve $\fS$ are not.
It is this gap that captures the singularity at $0$.

This observation leads to a refinement of the concept of embedding.
Let $X$ be a diffeological space and $\fS \subset X$ an embedded subset \cite[\textsection\,2.13]{PIZ13}.

Let us begin by defining two groupoids:

\begin{itemize}
  \item[(A)] Let $\KK_{\fS}$ be the groupoid of germs of local diffeomorphisms of $\fS$:
  $$
    \left\{
    \begin{aligned}
    \Obj(\KK_{\fS})  & = \fS \\
    \Mor_{\fS}(x,x') & = \{\germ(\pphi)_x \mid \pphi \in \Diff_\loc(\fS), \ \pphi(x) = x'\}
    \end{aligned}
    \right.
    $$
  \item[(B)] Let $\KK_{X,\fS}$ be the groupoid of germs of local diffeomorphisms of $X$ that preserve $\fS$:
  $$
    \left\{
    \begin{aligned}
    \Obj(\KK_{X,\fS})  & = \fS \\
    \Mor_{X,\fS}(x,x') & = \{\germ(\pphi)_x \mid \pphi \in \Diff_\loc(X,\fS), \ \pphi(x) = x'\}
    \end{aligned}
    \right.
    $$
\end{itemize}

There is a natural restriction morphism $\Phi$ from $\KK_{X,\fS}$ to $\KK_{\fS}$:
$$
  \Phi_{\Obj} = \Id_{\fS}
  \quad \text{and} \quad
  \Phi_{\Mor} : \germ(\pphi)_x \mapsto \germ(\pphi\restriction \fS)_x.
  $$
This leads to the following definition:

\begin{Article}{Definition 1.}{}
  We shall say that an embedded subset $\fS \subset X$ is \emph{smoothly embedded} if the morphism $\Phi$,
  from $\KK_{X,\fS}$ to $\KK_{\fS}$, is surjective on arrows.
  That is,
  if $\Phi$ is a full functor.
\end{Article}

In other words,
the embedded subset $\fS$ is smoothly embedded if it is embedded and the germ of any local diffeomorphism of $\fS$ can be extended to a local diffeomorphism of $X$;
or, equivalently, if the germ of any local diffeomorphism of $\fS$ is the imprint
(or trace)
of a local diffeomorphism of $X$.

Thus,
the semi-cubic is indeed an embedded submanifold,
but it is not smoothly embedded.

This notion of embedded subsets naturally leads to a definition of smooth embeddings.

\begin{Article}{Definition 2.}{}
  Let $j : \Sigma \to X$ be an embedding.
  We shall say that $j$ is a \emph{smooth embedding} if $\fS = j(\Sigma)$ is smoothly embedded in $X$.
\end{Article}

\textbf{Conclusion:} The injection $j : t \mapsto (t^2, t^3)$ is an embedding,
but not a smooth embedding.

\vfill
\centerline{
\includegraphics[width=\textwidth]{Figures/Notes-semi-cubic.pdf}}
